\documentclass[../../main.tex]{subfiles}
\begin{document}

\begin{center}
  {\fontsize{48pt}{58pt}\selectfont \textbf{1964}\normalsize\textbf{年度}}
\end{center}

\begin{table}[htb]
  \centering
  \caption{出題テーマと難易度}
  \begin{tabular}{|c|p{9cm}|c|}
    \hline
    \textbf{問題番号} & \textbf{テーマ} & \textbf{難易度} \\
    \hline
    第1問 & 分数関数の最大最小 & \\
    \hline
    第2問 & 不等式（相加相乗平均） & \\
    \hline
    第3問 & 図形と内接円の数列和 & \\
    \hline
    第4問 & 微分（最小値・道のり） & \\
    \hline
    第5問 & 積分と大小比較 & \\
    \hline
    第6問 & 極限（区分求積・高次評価） & \\
    \hline
  \end{tabular}
\end{table}

\begin{table}[htb]
  \centering
  \caption{解答}
  \renewcommand{\arraystretch}{1.4}
  \begin{tabular}{|c|c|p{9cm}|}
    \hline
    \textbf{問題番号} & \textbf{小問} & \textbf{答え} \\
    \hline
    第1問 & (1) & ない \\
    \cline{2-3}
    & (2) & $\max w=\dfrac{z}{z-1}$ \\
    \cline{2-3}
    & (3) & $\max w=2$ \\
    \cline{2-3}
    & (4) & $k=\dfrac{3}{2}$ \\
    \hline
    第2問 & (1) & 証明問題 \\
    \cline{2-3}
    & (2) & $P\ge R\ge S\ge Q$ \\
    \hline
    第3問 & - & $T=\pi r^2\left(1+\sum_{\theta\in\{A,B,C\}}\dfrac{(1-\sin(\theta/2))^2}{4\sin(\theta/2)}\right)$；正三角形では$\dfrac{11\pi a^2}{96}$ \\
    \hline
    第4問 & - & $\angle QOA=\dfrac{\pi}{3}$ \\
    \hline
    第5問 & - & $G(t)<H(t)$ \\
    \hline
    第6問 & (1) & $-\dfrac{a+b+1}{2}$ \\
    \cline{2-3}
    & (2) & $3+a$ \\
    \hline
  \end{tabular}
\end{table}

\paragraph{試験形態}
6問．試験時間は要確認．

\end{document}
